% GENERATED FILE. Edit canonical lessons, then run npm run generate:books.
% canonical-source-hash: fnv1a64:1e14828d4ae3c10d
% canonical-lessons: KA-C463-seve, KA-C463-saulabhya, KA-C463-odagisu, KA-C463-gisar, KA-C463-resart

\chapter{Hotel Services}
\label{ch:ka-hotel-services}
\hlchaptermodality{463}

\begin{chapteropening}
\textbf{By the end of this chapter:} \emph{I can say a service, a facility, to provide, a water heater and a resort.}

Complete the last lesson of chapter 463: I can say a service, a facility, to provide, a water heater and a resort.
\end{chapteropening}

\section[\texorpdfstring{\kn{ಸೇವೆ} (sēve)}{ಸೇವೆ (sēve)}]{\kn{ಸೇವೆ} (sēve) — a service}

\label{lesson:KA-C463-seve}



\begin{quote}
Before the new one: say the Kannada for a heater, then the Kannada for to creak.
\end{quote}

\subsection*{You'll want to know: \kn{ಸೇವೆ}}
\textbf{\kn{ಸೇವೆ}} — \emph{sēve} — \textquotedblleft{}a service\textquotedblright{}.

\begin{quote}
\kn{ಈ} \kn{ಹೋಟೆಲ್ನ} \kn{ಸೇವೆ} \kn{ತುಂಬಾ} \kn{ಚೆನ್ನಾಗಿದೆ}.
\end{quote}

\begin{quote}
\emph{ī hōṭelna sēve tumbā cennāgide.}
\end{quote}

\begin{quote}
The service in this hotel is very good.
\end{quote}

\kn{ಕೊಠಡಿ} \kn{ಸೇವೆ}, room service; \kn{ಬಸ್} \kn{ಸೇವೆ}, a bus service.

\begin{cousinweb}
\textbf{\kn{ಸೇವೆ}} (\emph{sēve}) is borrowed from Sanskrit \emph{sēvā}, \textquotedblleft{}attendance on, service\textquotedblright{} (CDIAL 13595); Kannada ends it in -\emph{e}.
\end{cousinweb}

\begin{grammarlens}[title={What you've built}]
One more word for getting around.
\end{grammarlens}

\subsection*{Your turn}
\begin{itemize}
\raggedright
  \item \emph{Say it:} \emph{sēve}
  \item \emph{Say it:} \emph{sēve}, once more
  \item \emph{Say it:} \emph{ī hōṭelna sēve tumbā cennāgide.} — The service in this hotel is very good.
  \item \emph{Recall:} say the Kannada for a heater, then the Kannada for to creak, then say \emph{sēve} again
\end{itemize}

\begin{tcolorbox}[breakable,colback=teal!4,colframe=teal!35!black,title={Before you move on}]
What does \kn{ಸೇವೆ} mean? (\textquotedblleft{}A service\textquotedblright{}.) Say it once more.
\end{tcolorbox}
\section[\texorpdfstring{\kn{ಸೌಲಭ್ಯ} (saulabhya)}{ಸೌಲಭ್ಯ (saulabhya)}]{\kn{ಸೌಲಭ್ಯ} (saulabhya) — a facility; an amenity}

\label{lesson:KA-C463-saulabhya}



\begin{quote}
Before the new one: say the Kannada for to creak, then the Kannada for a service.
\end{quote}

\subsection*{You'll want to know: \kn{ಸೌಲಭ್ಯ}}
\textbf{\kn{ಸೌಲಭ್ಯ}} — \emph{saulabhya} — \textquotedblleft{}a facility; an amenity\textquotedblright{}.

\begin{quote}
\kn{ಲಾಡ್ಜ್ನಲ್ಲಿ} \kn{ಇಂಟರ್ನೆಟ್} \kn{ಸೌಲಭ್ಯ} \kn{ಇಲ್ಲ}.
\end{quote}

\begin{quote}
\emph{lāḍjnalli iṇṭarneṭ saulabhya illa.}
\end{quote}

\begin{quote}
The lodge has no internet facility.
\end{quote}

Plural: \kn{ಸೌಲಭ್ಯಗಳು}, \emph{saulabhyagaḷu}.

\kn{ಸೌಲಭ್ಯಗಳು}, facilities: water, power, parking.

\begin{cousinweb}
\textbf{\kn{ಸೌಲಭ್ಯ}} (\emph{saulabhya}) is borrowed from Sanskrit \emph{saulabhya}, \textquotedblleft{}easiness of attainment\textquotedblright{} (Monier-Williams); in Kannada it names what makes a stay easy.
\end{cousinweb}

\begin{grammarlens}[title={What you've built}]
One more word for getting around.
\end{grammarlens}

\subsection*{Your turn}
\begin{itemize}
\raggedright
  \item \emph{Say it:} \emph{saulabhya}
  \item \emph{Say it:} \emph{saulabhya}, once more
  \item \emph{Say it:} \emph{lāḍjnalli iṇṭarneṭ saulabhya illa.} — The lodge has no internet facility.
  \item \emph{Recall:} say the Kannada for to creak, then the Kannada for a service, then say \emph{saulabhya} again
\end{itemize}

\begin{tcolorbox}[breakable,colback=teal!4,colframe=teal!35!black,title={Before you move on}]
What does \kn{ಸೌಲಭ್ಯ} mean? (\textquotedblleft{}A facility; an amenity\textquotedblright{}.) Say it once more.
\end{tcolorbox}
\section[\texorpdfstring{\kn{ಒದಗಿಸು} (odagisu)}{ಒದಗಿಸು (odagisu)}]{\kn{ಒದಗಿಸು} (odagisu) — to provide; to supply}

\label{lesson:KA-C463-odagisu}



\begin{quote}
Before the new one: say the Kannada for a service, then the Kannada for a facility.
\end{quote}

\subsection*{You'll want to know: \kn{ಒದಗಿಸು}}
\textbf{\kn{ಒದಗಿಸು}} — \emph{odagisu} — \textquotedblleft{}to provide; to supply\textquotedblright{}.

\begin{quote}
\kn{ಹೋಟೆಲ್} \kn{ಪ್ರವಾಸಿಗರಿಗೆ} \kn{ಉಚಿತ} \kn{ತಿಂಡಿ} \kn{ಒದಗಿಸುತ್ತದೆ}.
\end{quote}

\begin{quote}
\emph{hōṭel pravāsigarige ucita tiṇḍi odagisuttade.}
\end{quote}

\begin{quote}
The hotel provides tourists with a free breakfast snack.
\end{quote}

Past (I …): \kn{ಒದಗಿಸಿದೆ}, \emph{odagiside}.

What is given takes -\kn{ಅನ್ನು} or no ending; the receiver takes -\kn{ಗೆ}.

\begin{cousinweb}
\textbf{\kn{ಒದಗಿಸು}} (\emph{odagisu}) is native Dravidian: the causative of \emph{odagu}, \textquotedblleft{}to come to hand, to be of use\textquotedblright{} (DEDR 609): to make a thing be at hand.
\end{cousinweb}

\begin{grammarlens}[title={What you've built}]
One more word for getting around.
\end{grammarlens}

\subsection*{Your turn}
\begin{itemize}
\raggedright
  \item \emph{Say it:} \emph{odagisu}
  \item \emph{Say it:} \emph{odagisu}, once more
  \item \emph{Say it:} \emph{hōṭel pravāsigarige ucita tiṇḍi odagisuttade.} — The hotel provides tourists with a free breakfast snack.
  \item \emph{Recall:} say the Kannada for a service, then the Kannada for a facility, then say \emph{odagisu} again
\end{itemize}

\begin{tcolorbox}[breakable,colback=teal!4,colframe=teal!35!black,title={Before you move on}]
What does \kn{ಒದಗಿಸು} mean? (\textquotedblleft{}To provide; to supply\textquotedblright{}.) Say it once more.
\end{tcolorbox}
\section[\texorpdfstring{\kn{ಗೀಸರ್} (gīsar)}{ಗೀಸರ್ (gīsar)}]{\kn{ಗೀಸರ್} (gīsar) — a water heater}

\label{lesson:KA-C463-gisar}



\begin{quote}
Before the new one: say the Kannada for a facility, then the Kannada for to provide.
\end{quote}

\subsection*{You'll want to know: \kn{ಗೀಸರ್}}
\textbf{\kn{ಗೀಸರ್}} — \emph{gīsar} — \textquotedblleft{}a water heater\textquotedblright{}.

\begin{quote}
\kn{ಬೆಳಗ್ಗೆ} \kn{ಗೀಸರ್} \kn{ಕೆಲಸ} \kn{ಮಾಡಲಿಲ್ಲ}.
\end{quote}

\begin{quote}
\emph{beḷagge gīsar kelasa māḍalilla.}
\end{quote}

\begin{quote}
In the morning the water heater did not work.
\end{quote}

In Indian English a geyser heats bath water.

\begin{cousinweb}
\textbf{\kn{ಗೀಸರ್}} (\emph{gīsar}) is a loan from English \emph{geyser}, which in Indian English names a bathroom water heater.
\end{cousinweb}

\begin{grammarlens}[title={What you've built}]
One more word for getting around.
\end{grammarlens}

\subsection*{Your turn}
\begin{itemize}
\raggedright
  \item \emph{Say it:} \emph{gīsar}
  \item \emph{Say it:} \emph{gīsar}, once more
  \item \emph{Say it:} \emph{beḷagge gīsar kelasa māḍalilla.} — In the morning the water heater did not work.
  \item \emph{Recall:} say the Kannada for a facility, then the Kannada for to provide, then say \emph{gīsar} again
\end{itemize}

\begin{tcolorbox}[breakable,colback=teal!4,colframe=teal!35!black,title={Before you move on}]
What does \kn{ಗೀಸರ್} mean? (\textquotedblleft{}A water heater\textquotedblright{}.) Say it once more.
\end{tcolorbox}
\section[\texorpdfstring{\kn{ರೆಸಾರ್ಟ್} (resārṭ)}{ರೆಸಾರ್ಟ್ (resārṭ)}]{\kn{ರೆಸಾರ್ಟ್} (resārṭ) — a resort}

\label{lesson:KA-C463-resart}



\begin{quote}
Before the new one: say the Kannada for to provide, then the Kannada for a water heater.
\end{quote}

\subsection*{You'll want to know: \kn{ರೆಸಾರ್ಟ್}}
\textbf{\kn{ರೆಸಾರ್ಟ್}} — \emph{resārṭ} — \textquotedblleft{}a resort\textquotedblright{}.

\begin{quote}
\kn{ನಾವು} \kn{ಕಾಡಿನ} \kn{ಹತ್ತಿರದ} \kn{ಒಂದು} \kn{ರೆಸಾರ್ಟ್ನಲ್ಲಿ} \kn{ತಂಗಿದೆವು}.
\end{quote}

\begin{quote}
\emph{nāvu kāḍina hattirada ondu resārṭnalli taṅgidevu.}
\end{quote}

\begin{quote}
We stayed at a resort near the forest.
\end{quote}

A place to stay for a holiday, often out of town.

\begin{cousinweb}
\textbf{\kn{ರೆಸಾರ್ಟ್}} (\emph{resārṭ}) is a loan from English \emph{resort}.
\end{cousinweb}

\begin{grammarlens}[title={What you've built}]
One more word for getting around.
\end{grammarlens}

\subsection*{Your turn}
\begin{itemize}
\raggedright
  \item \emph{Say it:} \emph{resārṭ}
  \item \emph{Say it:} \emph{resārṭ}, once more
  \item \emph{Say it:} \emph{nāvu kāḍina hattirada ondu resārṭnalli taṅgidevu.} — We stayed at a resort near the forest.
  \item \emph{Recall:} say the Kannada for to provide, then the Kannada for a water heater, then say \emph{resārṭ} again
\end{itemize}

\begin{tcolorbox}[breakable,colback=teal!4,colframe=teal!35!black,title={Before you move on}]
What does \kn{ರೆಸಾರ್ಟ್} mean? (\textquotedblleft{}A resort\textquotedblright{}.) Say it once more.
\end{tcolorbox}
